\unitopen{36}
\begin{candos}
\item przejść test końcowy B1 z minimum 80\% poprawnych odpowiedzi;
\item znać format egzaminu B1 i strategię na każdy z czterech modułów;
\item poprowadzić rozmowę o swojej pracy i planach przez pięć minut bez notatek;
\item napisać formalny e-mail i krótki tekst argumentacyjny na poziomie B1.
\end{candos}

\section{Słownictwo}
\vocabtable{36}

\section{Format egzaminu B1}
Egzaminy typu Goethe-Zertifikat B1 i telc Deutsch B1 mają cztery moduły: \de{Lesen}, \de{Hören}, \de{Schreiben}, \de{Sprechen}.
Dokładny przebieg, liczba zadań i czas pracy zmieniają się; \emph{zawsze} pobierz aktualny arkusz modelowy od wystawcy egzaminu
i przepracuj go w warunkach egzaminu. Przy Goethe-Zertifikat B1 część mówiona obejmuje m.in.\ wspólne zaplanowanie czegoś,
krótką prezentację i rozmowę na temat \verify{(sprawdź aktualny opis na stronie Goethe-Institut)}.

\begin{center}\small
\begin{tabular}{@{}p{0.2\linewidth}p{0.72\linewidth}@{}}
\toprule
\textsf{\textbf{Moduł}} & \textsf{\textbf{Strategia}}\\
\midrule
\de{Lesen} & najpierw pytania, potem tekst; szukaj słów kluczowych, nie tłumacz wszystkiego\\
\de{Hören} & przeczytaj zadania przed nagraniem, notuj liczby i imiona; nagranie często jest powtarzane\\
\de{Schreiben} & trzy zdania: cel, treść, zakończenie; stosuj stałe zwroty z jednostki 22\\
\de{Sprechen} & mów w pełnych zdaniach, prosisz o powtórzenie bez wstydu, reaguj na partnera\\
\bottomrule
\end{tabular}
\end{center}

\begin{zasada}[Reguła 80\%]
Jeśli w teście końcowym B1 masz mniej niż 80\% (czyli mniej niż 10 z 12), wróć do tygodni 25, 28 i 29, zanim przejdziesz do B2.
Przed B2 zrób też jeden pełny arkusz egzaminu modelowego B1, w czasie, bez słownika.
\end{zasada}

\begin{falle}
Egzamin sprawdza komunikację, nie perfekcję. Błąd strategiczny: zatrzymanie się na jednym trudnym słowie i utrata czasu.
Zasada: pomiń, wróć później. Drugi: uczenie się na pamięć gotowych wypowiedzi, które nie pasują do zadania. Trzeci: pisanie za długo kosztem
poprawności i spójności.
\end{falle}

\begin{szwajcaria}
W Szwajcarii dla spraw pobytowych i obywatelskich istnieje system testów \de{fide}, mierzący poziom A1--B1 z naciskiem na mówienie i
życie codzienne \verify{(wymagania zależą od rodzaju pozwolenia; obywatel UE zwykle nie musi mieć certyfikatu do pracy, sprawdź w SEM i kantonie)}.
Pracodawcy w IT w Zurychu rzadko wymagają formalnego certyfikatu na B2, ale poziom B2 realnie ułatwia pracę i rozmowę z klientem.
\end{szwajcaria}

\begin{praca}
Własny plan na B2: wybierz termin egzaminu, ustaw cel (wynik \de{bestanden}) i zaplanuj dwa pełne arkusze modelowe. W kalendarzu
zapisz: \de{Prüfungsanmeldung}, \de{Probeprüfung 1}, \de{Probeprüfung 2}.
\end{praca}

\section{Dialogi}
\textbf{Dialog 1 --- część mówiona (planowanie)}
\begin{dialog}
\turn{A}{Wir sollen gemeinsam eine Abschlussfeier für das Team planen. Was meinst du?}{Mamy wspólnie zaplanować uroczystość zespołu. Co o tym sądzisz?}
\turn{B}{Ich würde ein Restaurant in der Nähe vorschlagen, damit alle schnell dort sind.}{Zaproponowałbym restaurację w pobliżu, żeby wszyscy szybko tam dotarli.}
\turn{A}{Gute Idee. Aber wir müssen zuerst fragen, ob alle Zeit haben. Wann passt es dir?}{Dobry pomysł. Ale najpierw musimy zapytać, czy wszyscy mają czas. Kiedy ci pasuje?}
\turn{B}{Am Freitag nach der Arbeit wäre ideal, oder?}{W piątek po pracy byłoby idealnie, prawda?}
\end{dialog}

\section{Test końcowy B1}
Próg: 80\% (10 z 12).

\begin{cwiczenia}
\exitem{Połącz: \de{Ich lerne Deutsch.} (um \dots{} zu) \de{Ich arbeite in der Schweiz.}}{\de{Ich lerne Deutsch, um in der Schweiz zu arbeiten.}}
\exitem{Zdanie względne: \de{Das ist der Kollege. Ich arbeite mit ihm.}}{\de{Das ist der Kollege, mit dem ich arbeite.}}
\exitem{Passiv Präsens: \de{Man testet den Code.}}{\de{Der Code wird getestet.}}
\exitem{Konjunktiv II: \de{Wenn ich Zeit \dots{}, \dots{} ich kommen.}}{\de{Wenn ich Zeit hätte, würde ich kommen.}}
\exitem{Czasownik z przyimkiem: \de{Ich warte \dots{} den Bus.}}{\de{auf}}
\exitem{Pytanie z \de{wo(r)-}: \de{Ich denke an den Termin.}}{\de{Woran denkst du?}}
\exitem{Genitiv: \de{das Auto (der Bruder)}.}{\de{das Auto des Bruders}}
\exitem{Łącznik z inwersją: \de{Es ist teuer.} (trotzdem) \de{Wir kaufen es.}}{\de{Es ist teuer, trotzdem kaufen wir es.}}
\exitem{Plusquamperfekt: \de{Nachdem er \dots{} (ankommen, Plusquamperfekt), begann das Meeting.}}{\de{Nachdem er angekommen war, begann das Meeting.}}
\exitem{Uprzejma prośba o przesunięcie terminu (jedno zdanie).}{\de{Wäre es möglich, den Termin zu verschieben?}}
\exitem{Standup: zdanie o tym, co zrobiłeś wczoraj.}{np.\ \de{Gestern habe ich einen Bug behoben.}}
\exitem{Przetłumacz: „Chociaż byłem zmęczony, przyszedłem na spotkanie.”}{\de{Obwohl ich müde war, bin ich zum Meeting gekommen.}}
\end{cwiczenia}

\section{Trening słownictwa (generowany)}
\drill{36}

\weekend{\textbf{Zadanie końcowe B1.} Przeprowadź pięciominutową rozmowę z AI jako symulację egzaminu ustnego (planowanie, prezentacja,
rozmowa). Napisz formalny e-mail o 150 słowach i krótki tekst argumentacyjny. Zapisz wyniki i porównaj z A2.}

\begin{uwaga}[Tydzień buforowy B1]
Po tej jednostce zrób tydzień buforowy: powtórz karty Anki z całego B1, napisz jedną pracę pisemną i wykonaj pełny arkusz egzaminu
modelowego. Dopiero potem zaczynaj B2.
\end{uwaga}
